\section{Introduction}
\label{sec:introduction}

A Banach space has the \emph{fixed point property for nonexpansive maps}
if every nonexpansive selfmap of each nonempty closed bounded convex
subset has a fixed point. Here nonexpansiveness is measured in the given
norm. We prove the following theorem.

\begin{theorem}\label{thm:main}
Let $X$ be a real reflexive Banach space, and let $C\subseteq X$ be
nonempty, norm closed, bounded, and convex. If $F:C\to C$ satisfies
\[
  \norm{F(a)-F(b)}\le \norm{a-b}\qquad(a,b\in C),
\]
then $F$ has a fixed point.
\end{theorem}

Theorem~\ref{thm:main} establishes the positive assertion in the
reflexive-space fixed point problem for the given norm.

Bruck's theorem for commuting families gives a further consequence.

\begin{corollary}[Commuting families]\label{cor:commuting}
Let $X$ be a real reflexive Banach space, and let $C\subseteq X$ be
nonempty, norm closed, bounded, and convex. For any pairwise commuting
family $\mathcal S$ of nonexpansive selfmaps of $C$, the common fixed set
\[
  \operatorname{Fix}(\mathcal S)
  =\{x\in C:T(x)=x\text{ for every }T\in\mathcal S\}
\]
is nonempty and is a nonexpansive retract of $C$. That is, there is a
nonexpansive map $R:C\to\operatorname{Fix}(\mathcal S)$ that is the
identity on $\operatorname{Fix}(\mathcal S)$.
\end{corollary}

\begin{proof}
For the empty family, take $R=\operatorname{id}_C$. Otherwise,
Theorem~\ref{thm:main} applies to every nonempty closed bounded convex
subset $D$ of $C$, since a relatively norm-closed subset of the
norm-closed set $C$ is norm closed in $X$. Thus $C$
has Bruck's hereditary fixed point property and, by taking $D=C$, the
fixed point property. If $f:C\to C$ is nonexpansive and such a subset
$D$ is $f$-invariant, the same theorem applied to $f|_D$ gives a fixed
point in $D$. Hence $C$ also has Bruck's conditional fixed point
property. The ordinary fixed point and conditional fixed point properties
just proved, together with weak compactness of $C$ by reflexivity, satisfy
the corresponding hypotheses of \citet[Theorem~1]{Bruck1974}. That theorem
gives the claimed nonempty nonexpansive retract.
\end{proof}

The classical results of Browder and G\"ohde prove the fixed point
property for uniformly convex spaces
\cite{Browder1965,Gohde1965}. Kirk's theorem applies when the closed
bounded convex domain $C$ in a reflexive space has normal structure
\cite{Kirk1965}: every convex subset $D\subseteq C$ with more than one
point must contain a point $a$ with
$\sup_{z\in D}\norm{a-z}<\diam D$.
Kirk's proof uses the minimal invariant set reduction employed below,
with the normal-structure and center ideas of Brodski\u\i{} and
Milman credited there \cite[p.~1004]{Kirk1965}.
Karlovitz subsequently proved the fixed point property for a particular
norm on $\ell_2$ without normal structure \cite{Karlovitz1976}.
Thus these geometric conditions give sufficient criteria, but normal
structure is not necessary for the property.

Further results established the property for substantial classes beyond
the uniformly convex case. Maurey proved it for reflexive subspaces of
$L^1$ in the inherited norm \cite[Theorem~1]{Maurey1981}.
Lin proved the weak fixed point property for spaces with a
$1$-unconditional basis, meaning that changing signs of basis coefficients
preserves the norm \cite[Theorem~1]{Lin1985}. Here the weak fixed point
property restricts the domains to weakly compact convex sets; in a
reflexive space it is equivalent to the property considered here.
Garc\'ia-Falset, Llorens-Fuster, and Mazcu\~n\'an-Navarro proved the
property for uniformly nonsquare spaces \cite{GarciaFalset2006}; see also
the proof of \citet{Dowling2008}. This condition requires a
uniform gap below $1$ in at least one of
$\norm{(a+b)/2}$ and $\norm{(a-b)/2}$ for every pair $a,b$ in the unit ball.
These theorems concern the specified norm, with additional structure.

Two other results distinguish the hypotheses of
Theorem~\ref{thm:main}. Alspach constructed a weakly compact convex
subset of $L_1[0,1]$ with a fixed-point-free nonexpansive selfmap
\cite{Alspach1981}. Weak compactness of the domain in an arbitrary
Banach space is therefore insufficient. Dom\'inguez Benavides proved
that every reflexive Banach space admits an equivalent norm with the
fixed point property \cite[Corollary~1]{Benavides2009}. A change of norm
can change which maps are nonexpansive; the theorem here concerns the
original norm.

Hanebaly previously proposed a proof of this same original-norm
assertion \cite[Theorem~1]{Hanebaly2019}, and subsequently announced
a further proof \cite{Hanebaly2021}. More recent discussion of the
general question appears in \citet[Section~1]{Nourouzi2026}, who
records it as unresolved. These statements of the literature are distinct
from the argument given here. Our proof proceeds through an adaptive
anchor, compact families of predicted outputs, and weighted switching
on a tree.

\paragraph{The argument.}
Suppose a counterexample exists. The classical minimal invariant set
method reduces the problem to a separable weakly compact convex set
$K$ of diameter one. Approximate fixed points in $K$ become asymptotically
diametral: their distance from every prescribed point of $K$ tends to one.
This is the Goebel--Karlovitz lemma
\cite{Goebel1975,Karlovitz1976}, proved in
Lemma~\ref{lem:diametral} below.

The proof uses this diametral behavior through two constructions.
First, a compact semigroup of nonexpansive maps provides an anchor
$x\in K$. Finite convex combinations of increasingly accurate
approximate fixed point images of $x$ can return arbitrarily close
to $x$, even when the accuracy required at each choice depends on
all previous choices. Second, these images label an ordered tree.
At each node, a compact family contains the required weak limits of
ideal root outputs, with parameters for the convex coefficients still
to be chosen. These outputs model the state after a prescribed block
of leaves has been replaced by a common approximate fixed point.

Replacing the leaves from right to left gives a sequence of actual
root outputs. Their consecutive differences telescope, although
the maps defining the tree are nonlinear. Weighted estimates for
these differences produce, at every finite depth, one vector
detected by all the functionals along a branch. A summable error
budget controls the total weighted deficit without a lower bound
on the individual branch weights. Finally, weak compactness of the
unit ball of $\overline{\operatorname{span}}(K-K)$ gives one vector
detected along an infinite branch, while the same functionals tend
pointwise to zero on that space.

The adaptive choice of the anchor images, the compact families of
predicted outputs, and the weighted telescoping estimate are the
parts of the argument that go beyond the classical diametral
reduction. They are formulated separately so that their dependence
on convexity, nonexpansiveness, and compactness is explicit.
Section~\ref{sec:prelim} supplies the compactness and separation
facts used throughout, with their proofs in
Appendix~\ref{app:foundations}; Sections~\ref{sec:anchor} and~\ref{sec:tree}
construct the anchor and tree; Section~\ref{sec:switching} proves
the switching estimate and completes the contradiction.
